\clearpage
\VerdierPhysicalPageAnchor{0193}%
\begingroup\small\setlength{\parskip}{0pt}\setlength{\emergencystretch}{1.5em}
\setlength{\abovedisplayskip}{4pt plus 2pt minus 2pt}%
\setlength{\belowdisplayskip}{4pt plus 2pt minus 2pt}%

\VerdierSemanticAnchor{ch3d-2-2-6-proof-cont}
(2.1.1)。射 $p_0$ は高さ $1$ の初等射である。複体 $U_0^\bullet$ は次の形をもつ：

\VerdierSemanticAnchor{ch3d-2-2-6-u0-complex}
\[
\begin{aligned}
U_0^\bullet={}&0\longrightarrow0\longrightarrow\cdots\longrightarrow0
\longrightarrow L_n\longrightarrow L_{n-1}\longrightarrow\cdots\\
&\longrightarrow L_1\longrightarrow Z_0\longrightarrow L
\longrightarrow0\longrightarrow0\longrightarrow\cdots.
\end{aligned}
\]
性質 $(E_4)$ により、対象 $Z_0$ は $M'$ に属する。したがって、エピ射
$u_0\colon L'_0\longrightarrow Z_0$ が存在し、$L'_0\in M$ である。
(2.1.1) と同様に、$U_1^\bullet=U_0^\bullet(u_0,0)$ および
$p_1=p(u_0,0)\colon U_1^\bullet\longrightarrow U_0^\bullet$ とおく。

\VerdierSemanticAnchor{ch3d-2-2-6-iterated-elementary-morphisms}
このように続けると、帰納法により射の列
\[
p_i\colon U_i^\bullet\longrightarrow U_{i-1}^\bullet,
\qquad 1\leq i\leq n-1,
\]
\[
p_0\colon U_0^\bullet\longrightarrow V^\bullet,
\]
が得られる。ここで $p_i$ は高さ $-i+1$ の初等射であり、
$U_i^\bullet$ は次の形をもつ：

\VerdierSemanticAnchor{ch3d-2-2-6-ui-complex}
\[
\begin{aligned}
U_i^\bullet={}&0\longrightarrow0\longrightarrow\cdots
\longrightarrow L_n\longrightarrow L_{n-1}\longrightarrow\cdots
\longrightarrow L_{i+1}\longrightarrow Z_i\\
&\longrightarrow L'_{i-1}\longrightarrow\cdots\longrightarrow L'_0
\longrightarrow L\longrightarrow0\longrightarrow\cdots ,
\end{aligned}
\]
ただし $L'_j\in M$, $0\leq j\leq i-1$ である。したがって、これらの射
$p_i$ を合成すると、エピな擬同型

\VerdierSemanticAnchor{ch3d-2-2-6-composite}
\[
p\colon U_{n-1}^\bullet\longrightarrow V^\bullet.
\]
が得られる。ここで、複体

\VerdierSemanticAnchor{ch3d-2-2-6-unminus1-complex}
\[
\begin{aligned}
U_{n-1}^\bullet={}&0\longrightarrow0\longrightarrow\cdots
\longrightarrow L_n\longrightarrow Z_{n-1}\longrightarrow L'_{n-2}
\longrightarrow\cdots\\
&\longrightarrow L'_0\longrightarrow L\longrightarrow0
\longrightarrow\cdots ,
\end{aligned}
\]
において、対象 $Z_{n-1}$ は性質 $(E_4)$ により $M$ の二つの対象の拡大として
現れるので、$M$ に属する（性質 $(E_5)$ による）。次をおく：

\VerdierSemanticAnchor{ch3d-2-2-6-l-complexes}
\[
L^\bullet=0\longrightarrow0\longrightarrow L_n\longrightarrow L_{n-1}
\longrightarrow\cdots\longrightarrow L_0\longrightarrow0
\longrightarrow0\longrightarrow\cdots ,
\]
\[
\begin{aligned}
{L'}^\bullet={}&0\longrightarrow0\longrightarrow L_n
\longrightarrow Z_{n-1}\longrightarrow L'_{n-2}
\longrightarrow\cdots\\
&\longrightarrow L'_0\longrightarrow0\longrightarrow0
\longrightarrow\cdots ,
\end{aligned}
\]
また、$m\colon {L'}^\bullet\longrightarrow L^\bullet$ を $p$ が誘導する射とする。
$L^\bullet$（それぞれ ${L'}^\bullet$）の唯一の非零コホモロジー対象は $Y$ で次数
$0$ にある（それぞれ $L$ で次数 $0$ にある）。また $m$ がこれらの対象上に誘導する
射は $u$ である。射 $m$ の写像錐 $c(m)$ を考える：

\VerdierSemanticAnchor{ch3d-2-2-6-cone}
\[
\begin{aligned}
c(m)={}&\cdots\longrightarrow0\longrightarrow L_n
\longrightarrow L_n\oplus Z_{n-1}
\longrightarrow L_{n-1}\oplus L'_{n-2}\longrightarrow\cdots\\
&\longrightarrow L_1\oplus L'_0\longrightarrow L_0
\longrightarrow0\longrightarrow\cdots.
\end{aligned}
\]
\endgroup
